\documentclass[11pt,a4paper]{article}
\usepackage[T1]{fontenc}
\usepackage[utf8]{inputenc}
\usepackage[main=italian,provide=*]{babel}
\usepackage{amsmath,amssymb}
\usepackage{geometry}
\geometry{margin=2.3cm,top=2.5cm,bottom=2.7cm}
\usepackage{booktabs}
\usepackage[table,dvipsnames]{xcolor}
\usepackage{tcolorbox}
\tcbuselibrary{skins,breakable}
\usepackage{titlesec}
\usepackage{enumitem}
\usepackage{needspace}
\usepackage{fancyhdr}
\usepackage{graphicx}
\usepackage{hyperref}
\hypersetup{colorlinks=true,linkcolor=Sepia,citecolor=Sepia,urlcolor=Sepia}
\newcommand{\Tr}{\operatorname{Tr}}
\newcommand{\ket}[1]{\lvert #1 \rangle}
\newcommand{\bra}[1]{\langle #1 \rvert}

\definecolor{inkblue}{HTML}{1B3A5C}
\definecolor{lightblue}{HTML}{EAF1F8}
\definecolor{accent}{HTML}{B0562A}
\definecolor{softgray}{HTML}{6B6B6B}
\definecolor{okgreen}{HTML}{2E6B3E}
\definecolor{okgreenbg}{HTML}{EAF5EC}

\titleformat{\section}
  {\normalfont\Large\bfseries\color{inkblue}}
  {\thesection}{0.6em}{}[{\color{inkblue!40}\titlerule}]
\titlespacing{\section}{0pt}{0.9em}{0.4em}
\titleformat{\subsection}
  {\normalfont\large\bfseries\color{accent}}
  {}{0em}{}
\titlespacing{\subsection}{0pt}{0.6em}{0.2em}

\pagestyle{fancy}
\fancyhf{}
\renewcommand{\headrulewidth}{0.4pt}
\fancyhead[L]{\small\color{softgray}Il modello di rumore per il trimero ad anello}
\fancyhead[R]{\small\color{softgray}\thepage}
\fancyfoot[C]{\small\color{softgray}Tesi triennale, Samuele --- Relatore: Prof.\ Alessandro Chiesa}

\newtcolorbox{keyeq}[1][]{
  colback=lightblue, colframe=inkblue!70, boxrule=0.6pt,
  arc=2pt, left=8pt, right=8pt, top=3pt, bottom=3pt, #1
}
\newtcolorbox{okbox}[1]{
  colback=okgreenbg, colframe=okgreen!60, boxrule=0.6pt,
  arc=2pt, left=7pt, right=7pt, top=2pt, bottom=2pt,
  fonttitle=\bfseries\color{okgreen}, title={#1}
}

\title{\vspace{-1.5em}\color{inkblue}Il modello di rumore per il trimero ad anello\\[0.2em]
\Large\normalfont\color{softgray}Stesso modello del dimero, riusato su un circuito a 3 qubit}
\author{Note di lavoro --- Tesi triennale, Samuele\\ Relatore: Prof.\ Alessandro Chiesa}
\date{}

\begin{document}
\sloppy
\maketitle
\thispagestyle{fancy}
\vspace{-2.0em}

File di codice: \texttt{noise\_model\_trimero\_anello.py} (nuovo, wrapper
sottile su \texttt{noise\_model\_dimero.py} --- Sez.~\ref{sec:riuso}),
\texttt{validate\_noise\_model\_trimero\_anello.py} (nuovo, verifica
l'aggancio a 3 qubit). Derivazione
completa delle formule in \texttt{teoria\_modello\_rumore.tex} (comune al
dimero, non richiama nulla di specifico a $N=2$). Nessun punto di lavoro
specifico per il modello in sé: come per il dimero, è la base condivisa
da tutti gli stadi successivi della pipeline.

\tableofcontents

\section{Obiettivo}

Aprire la Parte 2 (rumore) sul trimero ad anello, dopo la chiusura della
Parte 2 sul dimero (relazione inviata, riscontro positivo,
\texttt{log\_decisioni.md}, 6 settembre 2026). Il relatore ha confermato
la topologia (anello, non catena, stessa data) e lo scope dei canali di
rumore era già stato fissato per il dimero
(\texttt{domande\_relatore.md}, sez.~12): \textbf{due soli canali},
errore di gate (depolarizzante, 1 e 2 qubit) ed errore di lettura
(readout, simmetrico come primo passo). $T_1/T_2$ esclusi, stessa
decisione esplicita.

Questo documento non deriva un nuovo modello: verifica che quello già
costruito e validato per il dimero si applichi correttamente a un
sistema a 3 spin, senza modifiche.

\section{Riuso del modulo: perché è diretto, non un'estensione}
\label{sec:riuso}

Il modulo \texttt{noise\_model\_dimero.py} costruisce il
\texttt{NoiseModel} con due chiamate:
\begin{keyeq}
\centering
\texttt{nm.add\_all\_qubit\_quantum\_error(err, [...])} \qquad
\texttt{nm.add\_all\_qubit\_readout\_error(ro\_err)}
\end{keyeq}
Entrambe applicano l'errore a \emph{tutti} i qubit del circuito su cui
il \texttt{NoiseModel} viene agganciato in fase di simulazione, quale
che sia il loro numero: non c'è alcun riferimento hard-coded a 2 qubit
nel modulo. Il riuso per l'anello (3 qubit) è quindi diretto, non
un'estensione da scrivere --- verificato empiricamente in
Sez.~\ref{sec:verifica}, non solo argomentato dalla lettura del codice.

\begin{okbox}{Nota di dipendenza, dichiarata esplicitamente}
\texttt{noise\_model\_trimero\_anello.py} è un \textbf{wrapper sottile}:
importa e ri-espone tal quale (\texttt{build\_noise\_model},
\texttt{eps\_to\_lambda}, \texttt{BASIS\_GATES} e le costanti di
calibrazione) da \texttt{noise\_model\_dimero.py}, senza ridefinire
alcuna logica. Stessa situazione, dichiarata con lo stesso principio, già
segnalata per \texttt{validate\_ansatz\_params\_dimero.py} in Parte 1 del
dimero: un file concettualmente di un'altra parte del progetto viene
riusato da un contesto diverso, e la dipendenza è dichiarata nel
docstring del modulo, non nascosta dietro un import silenzioso. Il
motivo del wrapper, invece di un import diretto da
\texttt{validate\_noise\_model\_trimero\_anello.py}: dare al trimero un
punto d'ingresso con lo stesso nome usato in ogni altra parte del
progetto (\texttt{noise\_model\_dimero.py}, e in futuro
\texttt{noise\_model\_trimero\_catena.py} se servirà), cosi' che il
pacchetto riproducibile del trimero elenchi un
\texttt{noise\_model\_trimero\_anello.py} accanto agli altri moduli senza
dover spiegare ogni volta l'assenza --- senza però duplicare la logica in
due copie da tenere sincronizzate.
\end{okbox}

Le formule del canale depolarizzante e dell'errore di lettura, e la
calibrazione di riferimento, sono identiche a quelle già derivate per il
dimero (\texttt{teoria\_modello\_rumore.tex},
\texttt{risultati\_modello\_rumore\_dimero.tex}) --- non ripetute qui per
esteso, solo richiamate dove servono per la verifica.

\begin{center}
\renewcommand{\arraystretch}{1.3}
\begin{tabular}{@{}llc@{}}
\toprule
\textbf{Canale} & \textbf{Gate di riferimento} & \textbf{Valore mediano}\\
\midrule
Errore di gate, 1 qubit & \texttt{sx} & $\varepsilon_{1q}=2.9\times10^{-4}$\\
Errore di gate, 2 qubit & \texttt{CZ} & $\varepsilon_{2q}=3.8\times10^{-3}$\\
Errore di lettura & --- & $p_\text{readout}=2.3\times10^{-2}$\\
\bottomrule
\end{tabular}
\end{center}
Fonte: J.~P.~T.~Stenger, G.~Bazargan, N.~T.~Bronn, D.~Gunlycke,
\emph{``Method for simulating open-system dynamics using mid-circuit
measurements on a quantum computer''}, arXiv:2504.15187 (2025) ---
Appendice B, mediane pubblicate per \texttt{ibm\_torino}.

\begin{figure}[h]
\centering
\includegraphics[width=0.62\linewidth]{figure/fig_depolarizzazione_bloch.pdf}
\caption{\textit{Contrazione del vettore di Bloch, $1-\lambda$, in
funzione di $\lambda$ --- stesso contenuto del documento sul dimero: la
fisica del canale depolarizzante a 1 qubito non dipende dal numero di
spin del sistema più ampio in cui il qubit è inserito. \textbf{Nota
tecnica:} il PDF originale del Project, letto in questa sessione, risulta
un archivio zip con anteprima JPEG (limite della piattaforma per i file
binari montati in \texttt{/mnt/project} in questo ambiente --- verificato
su più PDF, non isolato a questo file --- non un difetto del file
sorgente). Rigenerata qui come vero PDF vettoriale
(\texttt{generate\_figure\_depolarizzazione\_bloch.py}, stesso stile
grafico del progetto via \texttt{stile.py}) a partire dalla formula nota
($1-\lambda$) e dalle etichette/range esatti recuperati dal testo
estratto nel manifest dello stesso pacchetto zip, non inventati.}}
\end{figure}

\section{Verifica}
\label{sec:verifica}

Punto di lavoro usato per la verifica: VQE+DM dell'anello, ansatz
$W$-2q.6, $J=1,\,J'=0.4,\,b=b_c=2.4,\,D=0.15$ (punto confermato, mode
``B''). Parametri VQE riottimizzati in questa sessione (multistart
COBYLA + polish L-BFGS-B, stesso schema del dimero):
$E_\text{VQE}=-5.53437001739340$, $E_\text{esatto}=-5.53437001739350$,
fedeltà $=0.99999999999961$ --- coincide con la fidelity $\approx1$ già
registrata in Parte 1 (Fase 5, ansatz $W$-2q.6).

\begin{okbox}{Controllo 1 --- limite di rumore nullo}
$\varepsilon_{1q}=\varepsilon_{2q}=p_\text{readout}=0$: l'energia
calcolata su \texttt{AerSimulator(method="density\_matrix")} deve
riprodurre l'energia VQE di Parte 1.
\begin{center}
\renewcommand{\arraystretch}{1.25}
\begin{tabular}{@{}lc@{}}
\toprule
\textbf{Quantità} & \textbf{Valore}\\
\midrule
$E_\text{VQE}$ (registrato, Parte 1) & $-5.5343700174$\\
$E_\text{VQE}$ (density\_matrix, rumore$=0$) & $-5.5343700174$\\
$|$differenza$|$ & $2.66\times10^{-15}$\\
CNOT nel circuito transpilato & $6$\\
\bottomrule
\end{tabular}
\end{center}
\textbf{OK}: differenza a precisione di macchina, ben sotto la soglia
$10^{-8}$ richiesta.
\end{okbox}

\begin{okbox}{Controllo 2 --- NoiseModel con valori di riferimento (\texttt{ibm\_torino})}
\begin{center}
\renewcommand{\arraystretch}{1.25}
\begin{tabular}{@{}lc@{}}
\toprule
\textbf{Quantità} & \textbf{Valore}\\
\midrule
$E_\text{VQE}$ (rumore di riferimento) & $-5.37687489$\\
Scostamento da $E_\text{VQE}$ ideale & $+0.157495$\\
CNOT nel circuito transpilato & $6$\\
\bottomrule
\end{tabular}
\end{center}
\textbf{OK}: energia rumorosa più alta del fondamentale, come atteso
fisicamente per un canale che degrada verso lo stato massimamente misto.
\end{okbox}

\begin{okbox}{Controllo 3 --- scan rapido su $\varepsilon_{2q}$}
A parità di $\varepsilon_{1q}=2.9\times10^{-4}$,
$p_\text{readout}=2.3\times10^{-2}$:
\begin{center}
\renewcommand{\arraystretch}{1.25}
\begin{tabular}{@{}cccc@{}}
\toprule
$\varepsilon_{2q}$ & $E_\text{VQE}$ rumoroso & scostamento & CNOT\\
\midrule
$0$ & $-5.49600915$ & $+0.038361$ & $6$\\
$1.0\times10^{-3}$ & $-5.46445616$ & $+0.069914$ & $6$\\
$3.8\times10^{-3}$ & $-5.37687489$ & $+0.157495$ & $6$\\
$8.0\times10^{-3}$ & $-5.24760282$ & $+0.286767$ & $6$\\
$1.5\times10^{-2}$ & $-5.03765731$ & $+0.496713$ & $6$\\
\bottomrule
\end{tabular}
\end{center}
\textbf{OK}: scostamento monotono crescente in $\varepsilon_{2q}$, come
atteso. Il conteggio di CNOT resta $6$ per ogni valore --- stesso
pattern già osservato sul dimero (il conteggio dopo transpilazione
dipende dalla struttura del circuito, non dai valori numerici dei
parametri del rumore né, qui, da quelli dell'ansatz).
\end{okbox}

\section{Perché 6 CNOT, e perché è indipendente da $\varepsilon_{2q}$}

Il circuito VQE+DM dell'anello (ansatz $W$-2q.6) contiene tre blocchi
$W$ (uno per legame: $(1,2)$, $(2,3)$, $(3,1)$), ciascuno costruito come
$\texttt{cx}\cdot R_y(\theta)\cdot\texttt{cx}$ --- 2 CNOT per blocco,
$3\times2=6$ in totale, più le rotazioni a 1 qubit finali che non
aggiungono CNOT. La transpilazione (\texttt{optimization\_level=3}) non
trova ulteriori fusioni: a differenza del blocco di Trotter del dimero
(dove passi consecutivi identici possono fondersi), qui i tre blocchi
agiscono su coppie di qubit diverse e in sequenza fissa, quindi non c'è
struttura ripetuta da collassare. Il conteggio non dipende da
$\varepsilon_{2q}$ per lo stesso motivo per cui non dipendeva dai valori
di $\theta$ sul dimero: il \texttt{NoiseModel} agisce \emph{dopo} la
transpilazione, sul circuito già fissato --- cambia la probabilità di
errore per ciascun gate, non la struttura del circuito stesso.

\section{Cosa non entra in questo modello}

Identico al dimero, per coerenza esplicita: $T_1/T_2$ esclusi
(decisione del relatore); nessuna correlazione fra qubit (ogni canale
applicato indipendentemente qubit per qubit, o coppia per coppia per il
\texttt{cx}); readout simmetrico soltanto in questa fase base;
calibrazione a un solo punto nel tempo (mediane pubblicate).

\section{Considerazioni e cosa serve al passo successivo}

Il riuso diretto del modulo (Sez.~\ref{sec:riuso}) significa che lo
Stadio 2 (VQE con DM sotto rumore) non richiede nuova infrastruttura di
\emph{rumore}: il \texttt{NoiseModel} è già pronto,
\texttt{energia\_vqe\_su\_simulatore(noise\_model, params)} già scritta
in \texttt{validate\_noise\_model\_trimero\_anello.py} ne fissa lo schema
d'uso (transpilazione + simulatore a operatore densità). La funzione di
costo effettiva dello Stadio 2 non è però un riuso letterale di
quest'ultima: aggiunge il calcolo della fedeltà (assente qui, dove serve
solo l'energia per il controllo 1), quindi è una funzione nuova ma
strutturalmente analoga --- stesso schema del dimero, dove
\texttt{validate\_noise\_model\_dimero.py} e
\texttt{vqe\_dm\_rumoroso\_dimero.py} sono due funzioni distinte per lo
stesso motivo, non la stessa funzione riusata (verificato ora sul
codice, non assunto per analogia).

A differenza del dimero, però, la strategia di transpilazione a blocchi
per Trotter e correlatori (Stadi 3--4) va scritta da zero: verificato
esplicitamente che né \texttt{circuito\_correlazioni\_trimero\_
anello.py} né \texttt{trotter\_trimero\_anello.py} contengono
riferimenti a \texttt{BASIS\_GATES} --- non un'estensione di codice
esistente, un modulo nuovo da scrivere seguendo lo stesso principio già
fissato sul dimero (transpilare il singolo passo/blocco una volta,
comporre $N$ volte).

\begin{okbox}{Sintesi in una riga}
Stesso modello di rumore del dimero (due canali, calibrazione
\texttt{ibm\_torino}), riusato senza modifiche su un circuito a 3 qubit
--- verificato con tre controlli su ansatz $W$-2q.6, non assunto
dall'analogia con $N=2$.
\end{okbox}

\needspace{4\baselineskip}
\section*{File prodotti}
\addcontentsline{toc}{section}{File prodotti}

\texttt{noise\_model\_trimero\_anello.py} (wrapper sottile su
\texttt{noise\_model\_dimero.py}, Sez.~\ref{sec:riuso}),
\texttt{validate\_noise\_model\_trimero\_anello.py}, questo documento.

\end{document}
